\section{Proving Set Equations}

\begin{frame}
\centering
\Huge\textbf{Proving Set Equations}
\end{frame}

% --------------------------------------------------
% Question 1
% --------------------------------------------------

\begin{frame}{Question 1}
A data-processing system classifies records using two sets $A$ and $B$.

A developer claims that processing records belonging to both sets is equivalent to first finding all records in $A$ and then removing those outside $B$.

\textbf{Question:} Prove the set equation

$$
A\cap B=A-(A-B).
$$

\end{frame}

\begin{frame}{Answer 1}
To prove two sets are equal, we show that an arbitrary element $x$ belongs to the left-hand side exactly when it belongs to the right-hand side.

Let

$$
x\in A\cap B.
$$

Then

$$
x\in A
\quad\text{and}\quad
x\in B.
$$

Since $x\in A$ and $x\notin A-B$,

$$
x\in A-(A-B).
$$
Therefore,
$$
A\cap B\subseteq A-(A-B).
$$

\end{frame}

\begin{frame}{Answer 1}
Now assume
$$
x\in A-(A-B).
$$

Then
$$
x\in A
$$

and
$$
x\notin A-B.
$$

Since $x\in A$, the fact that $x\notin A-B$ means that $x\in B$.

Therefore,
$$
x\in A\cap B.
$$

Thus,
$$
A-(A-B)\subseteq A\cap B.
$$

Hence,
$$
\boxed{A\cap B=A-(A-B)}.
$$

\end{frame}

% --------------------------------------------------
% Question 2
% --------------------------------------------------

\begin{frame}{Question 2}
In a software system, $A$ represents users with administrator privileges and $B$ represents users who are currently active.

A security rule claims:

\textit{The set of active administrators is the same as the active users who are administrators.}

\textbf{Question:} Prove

$$
A\cap B=B\cap A.
$$

\end{frame}

\begin{frame}{Answer 2}
Let

$$
x\in A\cap B.
$$

Then

$$
x\in A
\quad\text{and}\quad
x\in B.
$$

Since the order of the conditions does not matter,
$$
x\in B
\quad\text{and}\quad
x\in A.
$$

Therefore,
$$
x\in B\cap A.
$$

Hence,
$$
A\cap B\subseteq B\cap A.
$$

\end{frame}

\begin{frame}{Answer 2}
Similarly, let

$$
x\in B\cap A.
$$

Then

$$
x\in B
\quad\text{and}\quad
x\in A.
$$

Therefore,
$$
x\in A\cap B.
$$

Thus,
$$
B\cap A\subseteq A\cap B.
$$

Since both inclusions hold,
$$
\boxed{A\cap B=B\cap A}.
$$

\end{frame}

% --------------------------------------------------
% Question 3
% --------------------------------------------------

\begin{frame}{Question 3}
A machine-learning system uses set $A$ for records selected by model 1 and set $B$ for records selected by model 2.

A developer claims that selecting records from either model and then selecting those common to both produces the same result as directly selecting the common records.

\textbf{Question:} Prove

$$
A\cap(A\cup B)=A.
$$

\end{frame}

\begin{frame}{Answer 3}
Let

$$
x\in A\cap(A\cup B).
$$

Then

$$
x\in A
$$

and

$$
x\in A\cup B.
$$

The first condition already establishes that
$$
x\in A.
$$

Therefore,
$$
A\cap(A\cup B)\subseteq A.
$$

\end{frame}

\begin{frame}{Answer 3}
Now let

$$
x\in A.
$$

Since $x\in A$, it follows that

$$
x\in A\cup B.
$$

Therefore,

$$
x\in A\cap(A\cup B).
$$

Hence,
$$
A\subseteq A\cap(A\cup B).
$$

Both inclusions hold, so
$$
\boxed{A\cap(A\cup B)=A}.
$$

\end{frame}

% --------------------------------------------------
% Question 4
% --------------------------------------------------

\begin{frame}{Question 4}
A database administrator removes records belonging to set $B$ from a collection $A$.

The administrator then claims that removing $B$ from $A$ is equivalent to keeping the records in $A$ that are not in $B$.

\textbf{Question:} Prove

$$
A-B=A\cap B^c.
$$

\end{frame}

\begin{frame}{Answer 4}
Let

$$
x\in A-B.
$$

By the definition of set difference,

$$
x\in A
\quad\text{and}\quad
x\notin B.
$$

Since $x\notin B$,
$$
x\in B^c.
$$

Therefore,
$$
x\in A\cap B^c.
$$

Thus,
$$
A-B\subseteq A\cap B^c.
$$

\end{frame}

\begin{frame}{Answer 4}
Now assume

$$
x\in A\cap B^c.
$$

Then
$$
x\in A
\quad\text{and}\quad
x\in B^c.
$$

Therefore,
$$
x\notin B.
$$

Hence,
$$
x\in A-B.
$$

Thus,
$$
A\cap B^c\subseteq A-B.
$$

Therefore,
$$
\boxed{A-B=A\cap B^c}.
$$

\end{frame}

% --------------------------------------------------
% Question 5
% --------------------------------------------------

\begin{frame}{Question 5}
A cybersecurity system uses $A$ for normal network traffic and $B$ for encrypted traffic.

A security analyst claims that the complement of traffic that is either normal or encrypted is exactly the traffic that is neither normal nor encrypted.

\textbf{Question:} Prove

$$
(A\cup B)^c=A^c\cap B^c.
$$

\end{frame}

\begin{frame}{Answer 5}
Let

$$
x\in(A\cup B)^c.
$$

Then

$$
x\notin A\cup B.
$$

Therefore, $x$ is neither in $A$ nor in $B$:

$$
x\notin A
\quad\text{and}\quad
x\notin B.
$$

Hence,
$$
x\in A^c
\quad\text{and}\quad
x\in B^c.
$$

Therefore,
$$
x\in A^c\cap B^c.
$$

\end{frame}

\begin{frame}{Answer 5}
Conversely, suppose

$$
x\in A^c\cap B^c.
$$

Then

$$
x\notin A
\quad\text{and}\quad
x\notin B.
$$

Therefore,

$$
x\notin A\cup B.
$$

Hence,
$$
x\in(A\cup B)^c.
$$

Thus,
$$
\boxed{(A\cup B)^c=A^c\cap B^c}.
$$

\end{frame}
